\chapter{Robot Chain Control}
\label{robot_chain_control}

\minitoc{}

\section{Introduction}
\label{rcc_introduction}

\section{Robot chain modelling}
\label{rcc_robot_chain_modelling}

The inertial frame, denoted as \lframe{I} when needed, is an earth-fixed, right-handed frame.
Throughout this work we adopt a local \ac{enu} convention unless otherwise stated: the
\vect{x}-axis points to the east, the \vect{y}-axis points to the north and the \vect{z}-axis
points up. The origin of this frame is located at the water surface (\cref{rcc_rcm_frames}).

The body-fixed frame, denoted as \lframe{B} when needed, is rigidly attached to the vehicle with
origin at the vehicle centre of gravity. It is a right-handed frame with the \vect{x}-axis pointing
forward, the \vect{y}-axis pointing to the left and the \vect{z}-axis pointing up
(\cref{rcc_rcm_frames}).

\begin{figure}
	\centering
	\inputtikz{figures/robot_modelling_frames}
	\caption{Schematic representation of the inertial and body-fixed frames.}
	\label{rcc_rcm_frames}
\end{figure}

\subsection{Single robot dynamics}
\label{rcc_rcm_single_robot_dynamics}

The modelling of a \acp{rov} is described by the classical six-degree-of-freedom rigid body model
formulated for marine craft, as detailed in~\cite{fossenHandbookMarineCraft2011}:
\begin{align}
	\vect{\dot{\eta}} & = \matr{T}\left(\vect{\eta}\right) \vect{\nu} \label{rcc_rcm_srd_kinematics} \\
	\lrdmatr{M}{rb} \vect{\dot{\nu}} + \lrdmatr{M}{a} \ldvect{\dot{\nu}}{r}
	                  & + \lrdmatr{C}{rb}\left(\vect{\nu}\right) \vect{\nu}
	+ \left( \lrdmatr{C}{a}\left(\ldvect{\nu}{r}\right) + \matr{D}\left(\ldvect{\nu}{r}\right) \right) \ldvect{\nu}{r}
	+ \vect{s}\left(\vect{\eta}\right) = \vect{\omega} + \vect{\tau}
	\label{rcc_rcm_srd_dynamics}
\end{align}

In these equations:
\begin{itemize}
	\item $\vect{\eta} = \left[x, y, z, \phi, \theta, \psi\right] \in \lframe{I}$ denotes the vehicle's position and orientation in the inertial frame.
	\item $\vect{\nu} = \left[u, v, w, p, q, r\right] \in \lframe{B}$ is the linear and angular velocity in the body-fixed frame.
	\item $\matr{T}\left(\vect{\eta}\right)$ is the transformation matrix mapping body frame to inertial frame.
	\item \ldvect{\nu}{r} represents the relative velocity between the vehicle and the surrounding water, accounting for current disturbances.
	\item \lrdmatr{M}{rb} and \lrdmatr{M}{a} are the rigid-body and added mass inertia matrices, respectively.
	\item \lrdmatr{C}{rb} and \lrdmatr{C}{a} are their associated Coriolis matrices.
	\item \matr{D} is the hydrodynamic damping matrix.
	\item $\vect{s}\left(\vect{\eta}\right)$ represents the hydrostatic restoring force.
	\item \vect{\omega} includes environmental disturbances such as wind and wave loads.
	\item \vect{\tau} is the control input.
\end{itemize}

Under the assumptions of irrotational underwater currents and operation without dynamic ballast
systems, the dynamic model simplifies to the following form~\cite{fossenHandbookMarineCraft2011}:
\begin{align}
	\matr{I} \dot{\vect{\nu}}
	+ \left( \matr{C}(\vect{\nu}) + \matr{D}(\vect{\nu}) \right) \vect{\nu}
	+ \vect{s}(\vect{\eta})
	= \vect{\omega} + \vect{\tau}
	\label{rcc_rcm_srd_dynamics_simple}
\end{align}

Here, \matr{I} is the generalized inertia matrix, \matr{C} is the generalized Coriolis matrix, and
\matr{D} is the generalized damping matrix.

Furthermore, \cref{sota_sm_um_kinematics,sota_sm_um_dynamics_simple} may be combined into a
state-space representation:
\begin{align}
	\begin{bmatrix}
		\lrdmatr{I}{d} & 0        \\
		0              & \matr{I}
	\end{bmatrix}
	\begin{bmatrix}
		\dot{\vect{\eta}} \\
		\dot{\vect{\nu}}
	\end{bmatrix} & =
	\begin{bmatrix}
		0 & \matr{T}(\vect{\eta})                        \\
		0 & -\matr{C}(\vect{\nu}) - \matr{D}(\vect{\nu})
	\end{bmatrix}
	\begin{bmatrix}
		\vect{\eta} \\
		\vect{\nu}
	\end{bmatrix}
	+
	\begin{bmatrix}
		0 \\
		\vect{\omega} + \vect{\tau} - \vect{s}(\vect{\eta})
	\end{bmatrix}
\end{align}

With \lrdmatr{I}{d} the identity matrix.

The transformation matrix $\matr{T}\left(\vect{\eta}\right)$ is the transformation matrix from the
world frame to the body frame:
\begin{align}
	\matr{T}\left(\vect{\eta}\right) & = \begin{bmatrix}
		                                     \lrmatr{R}{w}{r}\left(\eta\right) & \lrdmatr{0}{3\times3}     \\
		                                     \lrdmatr{0}{3\times3}             & \matr{J}\left(\eta\right)
	                                     \end{bmatrix}
\end{align}

with $\lrmatr{R}{w}{r}\left(\eta\right)$ is the rotation matrix that represents the orientation of
the robot in the inertial frame, $\matr{J}\left(\eta\right)$ is defined as:
\begin{align}
	\matr{J}\left(\eta\right) & =\begin{bmatrix}
		                             1 & \sin{\phi} \tan{\theta}   & \cos{\phi} \tan{\theta}   \\
		                             0 & \cos{\phi}                & -\sin{\phi}               \\
		                             0 & \sin{\phi} / \cos{\theta} & \cos{\phi} / \cos{\theta} \\
	                             \end{bmatrix}
\end{align}

\subsection{Cable modelling}
\label{rcc_rcm_cable_modelling}

\lipsum[1]

\subsection{Robot chain dynamics}
\label{rcc_rcm_robot_chain_dynamics}

\lipsum[1]

\section{Model predictive control}
\label{rcc_model_predictive_control}

\subsection{Definition}
\label{rcc_mpc_definition}

\lipsum[1]

\subsection{Tuning}
\label{rcc_mpc_tuning}

\lipsum[1]

\subsection{Simulation}
\label{rcc_mpc_simulation}

\lipsum[1]

\section{Path planning}
\label{rcc_path_planning}

\subsection{Definition}
\label{rcc_pp_definition}

\lipsum[1]

\subsection{Tuning}
\label{rcc_pp_tuning}

\lipsum[1]

\subsection{Simulation}
\label{rcc_pp_simulation}

\lipsum[1]

\section{Reactive control}
\label{rcc_reactive_control}

\subsection{Definition}
\label{rcc_rc_definition}

\lipsum[1]

\subsection{Tuning}
\label{rcc_rc_tuning}

\lipsum[1]

\subsection{Simulation}
\label{rcc_rc_simulation}

\lipsum[1]

\section{Conclusion}
\label{rcc_conclusion}